%!TEX root = ../hutbthesis_main.tex

\chapter{. 总结与展望}\label{ux603bux7ed3ux4e0eux5c55ux671b}

\section{总结}\label{ux603bux7ed3}

本文主要对基于运动模仿的生物合理肌肉骨骼运动控制算法进行系统的探究，针对目前肌肉骨骼运动模仿存在评价不规范、指标单一、实验难以再现等现象，创建起一套标准化、可量化、可重复的运动控制和评价体系。全文主要工作可以概括为以下几个方面：

第一，系统梳理了肌肉骨建模、运动模仿学习强化学习、生物力学约束、量化评价等领域的理论与技术，指出了当前领域重效果、轻评价、难复现的问题，确定了以标准化流程为研究方向的路线，并结合指标进行可复现在实验。

第二，把生物合理的肌肉骨骼运动控制系统整体设计搭建出来了。整个系统采用模块化分层结构，把数据处理、前向运动学、策略控制、物理仿真和量化评估这五个模块有机地整合到一起，协同工作。系统以轨迹模仿为最终目标，核心算法选用PPO强化学习。这种设计不仅让运动过程精准，还兼顾了生物的合理性，最终能实现流畅、稳定、贴合人体生理规律的运动控制效果。

第三，对运动数据进行处理的全过程设计与实现，即数据加载、去噪、标准化、对齐、增强等过程，提高数据的质量和策略的泛化能力。同时建立多维模仿损失和生物合理性约束机制，从关节角度、肌肉驱动、能量消耗、运动平滑性四个角度保证运动真实性、安全性。

第四，构建标准化考核流程，将MPJPE、帧覆盖率、任务成功率三项指标进行组合，由"视觉相近"提升为"量化一致"。利用固定试验集、随机种子、运行参数，确保实验结果能验证，为肌肉、骨体控制算法提供了科学可信的验证范式。

第五，系统实现和实验验证，证明本文所设计的控制算法及评价体系在行走、跨步、转身等日常动作上具有高精度、高流畅度、高鲁棒性，可以为康复机器人、虚拟人动画、外骨骼控制等领域提供量化评价和算法验证的基础。

\section{研究创新点}\label{ux7814ux7a76ux521bux65b0ux70b9}

建立可以重复的标准化评价体系，采用固定的流程、统一的指标、固化参数，解决了肌肉骨骼运动模仿研究无法复现、无法对比的行业难题，使实验结论更加客观可信。

采取多维度组合量化的方法，从运动精度、连续性、鲁棒性三个方面综合评估MPJPE、帧覆盖率和任务成功率三个指标，克服单个误差指标的缺陷，更好地反映控制策略的实际情况。

生物合理性与模仿精度深度融合，在控制和奖励中加入关节约束、肌肉力学、能耗、平滑性约束，使生成的运动既具有轨迹相似性，又具有人体生理结构和运动性使运动自然性、可用性得到提高。

轻量化、模块化、易扩展的整体架构模块化程度高、接口清楚、部署方便，可以直接用在后续算法迭代、动作扩展和硬件迁移上，具有较好的工程实用性。

\section{未来展望}\label{ux672aux6765ux5c55ux671b}

复杂的多样的动作库为下一步的研究指明了方向，可以在提高策略的适应性和鲁棒性上加入更多类型、错综复杂的环境，如坡过坎、路面坎坷、推挤的外部干扰、避障运动等等。

融合多模态生理信号加入肌电、惯性测量、地面反作用力这些真实的生理数据，达成“仿真-真实”的双向对标，从而加强运动的生物合理性及可信度。

提高自主决策和互动能力，将系统从被动模仿升级为主动、稳定、自适应的拟人运动控制，采用目标导向、意图预测、在线适应性纠偏等方法，更加贴近真实的人体运动智慧。

将算法迁移到多关节外骨、康复机器人、智能假肢等硬件上，针对实际设备的落地应用，进行模拟到现实的验证，使研究成果运用到康复医疗、助老助残、人机协同等多个领域。
